% title: Race Rating Score Maximum
% short-title: Race Rating Score Maximum
% abstract: We prove that the maximum possible sum of the final scores in the race-rating problem is n(n-1)/2. The proof uses a per-student averaging bound: a student with score s has average rank at most n-s. Summing this over all students and all races gives the upper bound, and a fixed ordering in every race attains equality.
% subjclass: 05A05, 05A20
% keywords: rankings, double counting, extremal combinatorics, Lean 4
% author: Mario Hernández M.

\subsection{Race Rating Score Maximum}

\begin{problemstatement}
Let \(n\) be a positive integer. A class of \(n\) students run \(n\) races, in
each of which they are ranked with no draws. A student is eligible for a rating
\((a, b)\) for positive integers \(a\) and \(b\) if they come in the top \(b\)
places in at least \(a\) of the races. Their final score is the maximum possible
value of \(a-b\) across all ratings for which they are eligible. Find the
maximum possible sum of all the scores of the \(n\) students.
\end{problemstatement}

\begin{theorem}
The maximum possible sum of the final scores is
\[
  \frac{n(n-1)}{2}.
\]
\end{theorem}

\begin{proof}
For a fixed student, let \(c_b\) be the number of races in which that student
finishes in one of the top \(b\) places. The best score of the student is
\[
  s=\max_{1\le b\le n}(c_b-b),
\]
since for a fixed \(b\) the largest eligible \(a\) is \(c_b\).

Choose a value of \(b\) for which \(s=c_b-b\). Then the student has at least
\(b+s\) races with rank at most \(b\). If ranks are written as
\(0,1,\ldots,n-1\), this means that at least \(b+s\) of the zero-based ranks
are at most \(b-1\), while all remaining ranks are at most \(n-1\). Therefore
the student's total zero-based rank sum is at most
\[
  (b+s)(b-1)+(n-b-s)(n-1)\le n(n-1)-ns.
\]
Equivalently,
\[
  ns+\text{(student's zero-based rank sum)}\le n(n-1).
\]

Summing this inequality over all students gives
\[
  n\sum s_i+\sum_i\sum_{\text{races}} r_{i,\text{race}}
  \le n^2(n-1).
\]
In each race the zero-based ranks are exactly \(0,1,\ldots,n-1\), so the total
rank sum over all \(n\) races is
\[
  n\cdot\frac{n(n-1)}2.
\]
Consequently
\[
  n\sum s_i\le n\cdot\frac{n(n-1)}2,
\]
and, since \(n>0\),
\[
  \sum s_i\le \frac{n(n-1)}2.
\]
This upper-bound step is the formal content of
\lean{Biblioteca.Demonstrations.raceRatingScoreSum_le_maximum}.

The bound is sharp. Let every race have the same finishing order. Then the
student who is always in position \(j\), with \(j=1,\ldots,n\), has score
\(n-j\): choosing \(b=j\) gives \(a=n\), hence \(a-b=n-j\), and no larger score
is possible. Thus the scores are
\[
  n-1,n-2,\ldots,1,0,
\]
whose sum is \(\frac{n(n-1)}2\). In Lean this fixed-order sum is recorded as
\lean{Biblioteca.Demonstrations.fixedOrderRaceScores_sum}, and the combined
existence and upper-bound statement is
\lean{Biblioteca.Demonstrations.raceRatingMaximum_value}.
\end{proof}
